\chapter{Conjugation, Transmission Chains, and Structural Equivalence}\label{ch:conjugation}

The previous chapters built the static apparatus of an idea algebra: the
monoid $\IdMon$ (Chapter~\ref{ch:monoid}), the resonance pairing $\rho$
(Chapter~\ref{ch:resonance}), the Aufhebung decomposition
$\sigma+\nu+\eta$ (Chapter~\ref{ch:aufhebung}) and the emergence cocycle
$\omega$ (Chapter~\ref{ch:emergence}). What is still missing is a
\emph{theory of how ideas move}. An idea is rarely the property of a
single carrier; it passes through minds, languages, institutions, and
historical periods, and it is reshaped by each. Idea Theory accommodates
this dynamics with three closely linked devices. First, the invertible
ideas form a group $\Ideas^{\times}$, and conjugation by an element of
$\Ideas^{\times}$ is an automorphism of the entire structure: it
preserves $\rho$, the Aufhebung primitives, the cocycle $\omega$, and the
grading. Second, finite sequences of carriers can be folded into
\emph{transmission chains}, whose accumulated emergence telescopes
through $\omega$. Third, the right notion of sameness for two idea
algebras --- \emph{structural equivalence} --- is a graded monoid
isomorphism that intertwines all of the primitives at once.

These three devices, taken together, are the algebraic backbone of every
later interpretive chapter, from Bourdieu's habitus to Kuhn's
revolutions. The Lean development distributes them across three theorem
files: conjugation in \texttt{Theorems7}, transmission chains in
\texttt{Theorems7}--\texttt{Theorems8}, and the universal property of
structural equivalence in \texttt{Theorems8}--\texttt{Theorems9}. We
recapitulate that material here in mathematical prose, with no Lean
syntax in the body.

\section{Invertible ideas as a group}\label{sec:invertibles}

Let $\IdeaAlg = (\Ideas,\compose,\unit,\rho,\sigma,\nu,\eta,\omega,\deg)$
be a fixed idea algebra. By Theorem~1 of Chapter~\ref{ch:monoid},
$(\Ideas,\compose,\unit)$ is a monoid. We isolate the elements that
behave best with respect to that monoid.

\begin{definition}[Invertible idea]\label{def:invertible}
An idea $g\in\Ideas$ is \emph{invertible} if there exists an idea
$g^{-1}\in\Ideas$ with $g\compose g^{-1} = g^{-1}\compose g = \unit$. The
collection of invertibles is denoted $\Ideas^{\times}$. We write
$\mathsf{e}\colon \Ideas^{\times}\hookrightarrow\Ideas$ for the
underlying inclusion.
\end{definition}

That $g^{-1}$ exists at all is a strong condition: the typical idea
(``the idea of triangularity'', ``the proof of Lagrange's theorem'')
admits no two-sided inverse with respect to composition. The invertible
ideas are the \emph{operators} or \emph{transformations} on the algebra
--- the cognitive analogues of conjugating actions, change-of-frame
maps, and (in cultural settings) institutional relabelings.

\begin{lemma}[Uniqueness of inverses]\label{lem:inv-unique}
If $g\compose h=\unit=h\compose g$ and $g\compose h'=\unit=h'\compose g$,
then $h=h'$.
\end{lemma}

\begin{proof}
By associativity, $h = h\compose\unit = h\compose(g\compose h') =
(h\compose g)\compose h' = \unit\compose h' = h'$. The two cancellations
use only the monoid axioms, so the lemma holds in the bare structure
$(\Ideas,\compose,\unit)$ as well; this is the form verified as
\texttt{IdeaTheory.ident\_unique\_left} together with the right-sided
analogue.
\end{proof}

\begin{proposition}[Invertibles form a group]\label{prop:inv-group}
The set $\Ideas^{\times}$ is closed under $\compose$, contains $\unit$,
and is closed under $g\mapsto g^{-1}$. With the operations inherited
from $\Ideas$ it is a group.
\end{proposition}

\begin{proof}
The unit is its own inverse, so $\unit\in\Ideas^{\times}$. If
$g,h\in\Ideas^{\times}$ then $(h^{-1}\compose g^{-1})\compose
(g\compose h)= h^{-1}\compose(g^{-1}\compose g)\compose h =
h^{-1}\compose h = \unit$, and the symmetric calculation gives the other
side, so $g\compose h$ is invertible with inverse $h^{-1}\compose
g^{-1}$. Closure under inversion is automatic from the symmetry of the
defining equations. Associativity, the unit law, and the inverse law
together yield the group axioms; this is essentially the argument of
\cite{Bourbaki1989} for the units of any monoid.
\end{proof}

\subsection*{The inner-automorphism action}

Once $\Ideas^{\times}$ is a group, it acts on $\Ideas$ in the standard
way.

\begin{definition}[Inner action / conjugation]\label{def:inner-action}
For $g\in\Ideas^{\times}$, define the \emph{conjugation map}
\[
  \conj{g}{(\,\cdot\,)} \colon \Ideas \to \Ideas,
  \qquad
  \conj{g}{a} \;\eqdef\; g\compose a\compose g^{-1}.
\]
We also write $g\cdot a$ for $\conj{g}{a}$ when no ambiguity arises.
\end{definition}

\begin{lemma}[Action axioms]\label{lem:inner-action}
The map $g\mapsto \conj{g}{(\,\cdot\,)}$ is a left action of
$\Ideas^{\times}$ on $\Ideas$: $\conj{\unit}{a} = a$ and
$\conj{(g\compose h)}{a} = \conj{g}{(\conj{h}{a})}$ for all
$g,h\in\Ideas^{\times}$ and all $a\in\Ideas$.
\end{lemma}

\begin{proof}
Both equalities are immediate from associativity and the inverse-of-a-product
formula of \Cref{prop:inv-group}: $\conj{(g\compose h)}{a} = (g\compose
h)\compose a\compose (g\compose h)^{-1} = g\compose(h\compose a\compose
h^{-1})\compose g^{-1} = \conj{g}{(\conj{h}{a})}$.
\end{proof}

The orbit of $a$ under this action is the set $\{\conj{g}{a} :
g\in\Ideas^{\times}\}$; we call it the \emph{conjugacy class} of $a$.
Two ideas in the same conjugacy class are, by definition, related by an
``inner change of viewpoint''. The headline of \S\ref{sec:conj-thm}
will be that they agree on every primitive of the algebra.

\subsection*{Powers and centralizers}

Two further constructions on $\Ideas^{\times}$ will be useful in
later chapters and are recorded here in their basic form.

\begin{definition}[Powers]\label{def:powers}
For $g\in\Ideas^{\times}$ and $n\in\ZZ$, define $g^{0}=\unit$,
$g^{n+1} = g\compose g^n$, and $g^{-n} = (g^{-1})^n$. The map $\ZZ\to
\Ideas^{\times}$, $n\mapsto g^n$, is a group homomorphism whose
image is the cyclic subgroup generated by $g$.
\end{definition}

\begin{lemma}[Powers and conjugation]\label{lem:pow-conj}
For all $g\in\Ideas^{\times}$, $a\in\Ideas$, and $n\in\ZZ$,
$\conj{g^n}{a} = \conj{g}{(\conj{g^{n-1}}{a})}$, and the iterated
conjugation acts on chains by
$\conj{g^n}{(\chain{a_1,\dots,a_k})}
= \chain{\conj{g^n}{a_1},\dots,\conj{g^n}{a_k}}$.
\end{lemma}

\begin{proof}
The first equation is the action axiom of \Cref{lem:inner-action}
applied iteratively. The second uses that $C_{g^n}$ is a monoid
homomorphism (\Cref{thm:conjugation}(1)) together with the
chain-image lemma (\Cref{lem:image-chain}); concretely, expand
$\chain{a_1,\dots,a_k}$ as the iterated $\compose$, push $C_{g^n}$
through each step, and re-assemble the resulting chain. The
verified analogue for the bare-monoid case follows from
\texttt{IdeaTheory.compPow\_succ} composed with
\texttt{IdeaTheory.evalWord\_mapWord}.
\end{proof}

\begin{definition}[Centralizer]\label{def:centralizer}
The \emph{centralizer} of $a\in\Ideas$ in $\Ideas^{\times}$ is
\[
  Z_{\Ideas^{\times}}(a) \;\eqdef\; \{\,g\in\Ideas^{\times}\,:\,
    \conj{g}{a} = a\,\}.
\]
\end{definition}

\begin{proposition}[Centralizer is a subgroup]\label{prop:cent-sub}
For every $a\in\Ideas$, $Z_{\Ideas^{\times}}(a)$ is a subgroup of
$\Ideas^{\times}$. Its conjugacy-class orbit on $a$ is the singleton
$\{a\}$ and its quotient $\Ideas^{\times}/Z_{\Ideas^{\times}}(a)$
parametrizes the conjugacy class of $a$.
\end{proposition}

\begin{proof}
Closure: if $g,h$ both fix $a$ under conjugation, then
$\conj{(g\compose h)}{a} = \conj{g}{(\conj{h}{a})} = \conj{g}{a} = a$.
The unit fixes $a$ since $\conj{\unit}{a} = a$. Inverse: if
$\conj{g}{a} = a$, then $a = \conj{g^{-1}}{(\conj{g}{a})} =
\conj{g^{-1}}{a}$. The orbit-stabilizer statement is the standard
theorem for any group action, here applied to $C$.
\end{proof}

\section{The conjugation theorem}\label{sec:conj-thm}

The point of conjugation is not merely that it acts; it is that the
action is by automorphisms of the \emph{full} idea-algebra structure,
not just of the bare monoid. We collect the components.

\begin{theorem}[Conjugation theorem]\label{thm:conjugation}
For every $g\in\Ideas^{\times}$ the map $C_g\colon a\mapsto
\conj{g}{a}$ is a bijection of $\Ideas$ that satisfies, for all
$a,b\in\Ideas$ and (where the corresponding primitive is defined) for
all relevant arguments,
\begin{enumerate}
  \item $C_g(\unit) = \unit$ and $C_g(a\compose b) = C_g(a)\compose C_g(b)$
        (monoid automorphism);
  \item $\rho(C_g(a),C_g(b)) = \rho(a,b)$ (resonance invariance);
  \item $\sigma(C_g(a),C_g(b)) = C_g(\sigma(a,b))$,
        $\nu(C_g(a),C_g(b)) = C_g(\nu(a,b))$,
        $\eta(C_g(a),C_g(b)) = C_g(\eta(a,b))$ (Aufhebung equivariance);
  \item $\omega(C_g(a),C_g(b)) = \omega(a,b)$ (cocycle invariance);
  \item $\deg(C_g(a)) = \deg(a)$ (grading invariance).
\end{enumerate}
Hence the assignment $g\mapsto C_g$ is a homomorphism from
$\Ideas^{\times}$ into the automorphism group of the full
idea-algebra structure on $\Ideas$.
\end{theorem}

\begin{proof}
\emph{Bijectivity.} The inverse of $C_g$ is $C_{g^{-1}}$ by
\Cref{lem:inner-action}.

\emph{Monoid automorphism.} $C_g(\unit) = g\compose\unit\compose g^{-1}
= g\compose g^{-1} = \unit$. For composition,
\[
C_g(a)\compose C_g(b) = (g\compose a\compose g^{-1})\compose
  (g\compose b\compose g^{-1}) = g\compose a\compose b\compose g^{-1}
  = C_g(a\compose b).
\]
This is the only place where invertibility is needed for the monoid
clause: the $g^{-1}\compose g = \unit$ in the middle is what allows the
two carriers $a$ and $b$ to fuse without creating extra terms. The
verified analogue is the homomorphism law of the structure-preserving
maps \texttt{IdeaTheory.IdeaHom}.

\emph{Resonance invariance.} Recall from Chapter~\ref{ch:resonance} that
$\rho$ extends bilinearly to the free $\RR$-vector space $\RR\Ideas$,
and that for any monoid endomorphism $f$ of $\Ideas$ the pullback
$f^{*}\rho(a,b)\eqdef \rho(f(a),f(b))$ is again a bilinear pairing. The
map $C_g$ is, by part~(1), a monoid automorphism. Resonance was
postulated to be invariant under \emph{any} monoid automorphism induced
by an invertible carrier (Chapter~\ref{ch:resonance}, Axiom~R3); apply
this to $C_g$. Concretely, since $g$ is invertible, conjugation acts on
$\RR\Ideas$ by an invertible linear map whose adjoint with respect to
$\rho$ is its inverse, and the desired equality
$\rho(C_g a, C_g b) = \rho(a,b)$ is the explicit formula for this
adjoint.

\emph{Aufhebung equivariance.} The Aufhebung decomposition
$a\compose b = \sigma(a,b) + \nu(a,b) + \eta(a,b)$ was characterized in
Chapter~\ref{ch:aufhebung} as the unique (modulo $\ker\rho$)
decomposition into a $\rho$-symmetric, a $\rho$-antisymmetric, and a
$\rho$-orthogonal piece. Each of these characterizations is preserved
by any $\rho$-isometric monoid automorphism, and $C_g$ is exactly such
by parts~(1) and~(2). Hence $C_g$ permutes the three components: the
$\sigma$-piece of $C_g(a)\compose C_g(b)$ is the $\sigma$-piece of
$C_g(a\compose b) = C_g(\sigma(a,b)+\nu(a,b)+\eta(a,b))$, which by
linearity is $C_g(\sigma(a,b))$, and similarly for $\nu$ and $\eta$.

\emph{Cocycle invariance.} The cocycle $\omega$ was defined as the
defect of any normalized valuation $v\colon\Ideas\to A$ from being a
homomorphism: $v(a\compose b) = v(a)+v(b)+\omega(a,b)$. For an
invertible $g$ we have $v(\unit) = v(g\compose g^{-1}) = v(g)+v(g^{-1})
+\omega(g,g^{-1})$, so $v(g^{-1}) = -v(g)-\omega(g,g^{-1})$. Compute
\[
v(C_g(a)) = v(g\compose a\compose g^{-1})
= v(g)+v(a)+\omega(g,a) + v(g^{-1}) + \omega(g\compose a, g^{-1}).
\]
Substituting and collecting using the cocycle identity
$\omega(g,a)+\omega(g\compose a,g^{-1}) =
\omega(g,a\compose g^{-1})+\omega(a,g^{-1})$ (this is precisely
$\partial\omega = 0$ specialized to the triple $(g,a,g^{-1})$), one
finds $v(C_g(a)) = v(a) + (\text{coboundary of } v_g)$. The coboundary
piece is killed when one passes to the cohomology class, but here we
have an even stronger statement: applying the same calculation to
$C_g(a)\compose C_g(b)$ and using parts~(1) and the cocycle identity
\emph{twice} yields $\omega(C_g(a), C_g(b)) = \omega(a,b)$ on the
nose. The detailed bookkeeping is recorded as
\texttt{IdeaTheory.theorem\_4\_*} in the Lean development.

\emph{Grading invariance.} By Chapter~\ref{ch:grading}, the grading
$\deg\colon\Ideas\to G$ is a monoid homomorphism into a totally ordered
abelian group. Since $G$ is abelian and $C_g$ is a monoid
automorphism,
\[
\deg(C_g(a)) = \deg(g) + \deg(a) + \deg(g^{-1})
            = \deg(g) + \deg(g^{-1}) + \deg(a)
            = \deg(\unit) + \deg(a) = \deg(a). \qedhere
\]
\end{proof}

\begin{corollary}[Conjugation acts by structure isomorphisms]
\label{cor:conj-iso}
For every $g\in\Ideas^{\times}$, the map $C_g$ is an isomorphism of
the full idea algebra $\IdeaAlg$ onto itself. The image of $g\mapsto
C_g$ lies in the kernel of the natural map from $\mathrm{Aut}(\IdeaAlg)$
to the permutation group of conjugacy classes, and is normal in
$\mathrm{Aut}(\IdeaAlg)$.
\end{corollary}

\begin{proof}
The first sentence is the conclusion of \Cref{thm:conjugation}.
Normality follows by the same argument as for inner automorphisms of an
abstract group: if $\Phi\in\mathrm{Aut}(\IdeaAlg)$ then
$\Phi\circ C_g \circ\Phi^{-1} = C_{\Phi(g)}$, and $\Phi(g)$ is again
invertible because $\Phi$ preserves $\unit$ and $\compose$.
\end{proof}

\begin{remark}[Why \emph{inner} matters]\label{rem:inner-meaning}
The conjugation theorem is the algebraic explication of an old
philosophical observation: an idea is the same idea under a change of
expressive frame, provided the change of frame is itself an
\emph{idea} (here, an invertible one). Outer transformations --- those
not realized by any element of $\Ideas$ --- can in principle alter
$\rho$, $\omega$, or $\deg$; only the inner ones are guaranteed not to.
This will be the formal substrate for Bourdieu's habitus
(\S\ref{zoo:bourdieu}) and Foucault's épistémè
(\S\ref{zoo:foucault}). One may also read the theorem as a
non-triviality result: were the conjugation action ever to act
trivially on $\rho$ or $\omega$ \emph{outside} of inner
automorphisms, the algebraic distinction between inner and outer
would collapse, and with it the philosophical work the theorem is
called upon to perform. The conjugation theorem secures the
distinction by exhibiting it as a structural feature, not a
definitional stipulation.
\end{remark}

\begin{remark}[Linearity of the action]\label{rem:linear-action}
Because $C_g$ is a monoid endomorphism that fixes $\unit$, it
extends uniquely to an $\RR$-linear endomorphism of the free
$\RR$-vector space $\RR\Ideas$ on $\Ideas$. By
\Cref{thm:conjugation}(2) this extension is an isometry of the
bilinear pairing $\rho$. The collection
$\{C_g : g\in\Ideas^{\times}\}$ thus sits inside the orthogonal
group $O(\RR\Ideas,\rho)$, and the failure of inner-conjugation to
exhaust this orthogonal group is itself an algebraic invariant of
$\IdeaAlg$. We will return to this point in Chapter~\ref{ch:cogsci}
when discussing learned representational symmetries.
\end{remark}

\section{Transmission chains}\label{sec:chains}

We now turn from single conjugating actions to the iterated dynamics of
\emph{transmission}. The basic combinatorial object is a finite
sequence of carriers, folded by composition.

\begin{definition}[Transmission chain]\label{def:chain}
For a finite list $(a_1,\dots,a_n)$ of ideas, with $n\ge 0$, the
\emph{transmission chain} is the right-fold
\[
  \chain{a_1,\dots,a_n}
  \;\eqdef\;
  a_1\compose(a_2\compose(\cdots\compose(a_n\compose\unit)\cdots))
  = \prod_{i=1}^{n} a_i \quad (\text{empty product } = \unit).
\]
This is the function defined recursively in the Lean development as
$\mathtt{chain}\,[\,]=\unit$ and $\mathtt{chain}\,(a::\mathit{xs})=
a\compose\mathtt{chain}\,\mathit{xs}$; see
\texttt{IdeaTheory.chain} in \texttt{Theorems7.lean}.
\end{definition}

The basic algebraic facts about chains are immediate:

\begin{lemma}[Chain calculus]\label{lem:chain-calc}
For all finite lists $\mathit{xs}, \mathit{ys}$ of ideas:
\begin{enumerate}
  \item $\chain{} = \unit$, $\chain{a} = a$, $\chain{a,b} = a\compose b$;
  \item $\chain{\mathit{xs}\!\frown\!\mathit{ys}}
        = \chain{\mathit{xs}}\compose\chain{\mathit{ys}}$
        \emph{(append)};
  \item inserting $\unit$ at any position in a chain leaves the chain
        invariant;
  \item the chain $\chain{a_1,\dots,a_n}$ is invariant under any
        rebracketing of the underlying product;
  \item if $a$ is self-fixed in the sense that $a\compose a = a$, then
        $\chain{\underbrace{a,\dots,a}_{n+1}} = a$ for every $n\ge 0$.
\end{enumerate}
\end{lemma}

\begin{proof}
(1) is by definition. (2) is by induction on $\mathit{xs}$, using
associativity in the inductive step --- this is
\texttt{IdeaTheory.chain\_append}. (3) reduces to (2) and the unit law.
(4) is the generalized associative law for monoids
\cite[ch.~II]{Bourbaki1989}. (5) is by induction: at $n=0$ the claim is
$\chain{a} = a$, and at $n+1$ we have $\chain{a,a,\dots,a} = a\compose
\chain{a,\dots,a} = a\compose a = a$. The verified version is
\texttt{IdeaTheory.selfFixedPoint\_chain\_replicate}.
\end{proof}

\subsection*{Telescoping under $\omega$ --- the chain rule}

The interesting content of a chain is not its composite (which is just
an element of $\Ideas$) but the \emph{accumulated emergence} along it.
Recall from Chapter~\ref{ch:emergence} that $\omega\colon\Ideas\times
\Ideas\to A$ is a normalized 2-cocycle and that for any normalized
valuation $v\colon\Ideas\to A$ one has $v(a\compose b) = v(a)+v(b)+
\omega(a,b)$. Iterating gives a closed-form telescoping.

\begin{theorem}[Chain rule for $\omega$]\label{thm:chain-rule}
Let $v\colon\Ideas\to A$ be a normalized valuation with cocycle
$\omega$. For any chain $\chain{a_1,\dots,a_n}$ with $n\ge 1$,
\[
  v\!\bigl(\chain{a_1,\dots,a_n}\bigr)
  \;=\;
  \sum_{i=1}^{n} v(a_i)
  \;+\;
  \sum_{k=1}^{n-1}
    \omega\!\bigl(a_k,\, \chain{a_{k+1},\dots,a_n}\bigr).
\]
The terms in the second sum are the \emph{telescoped emergence
defects}; they vanish identically in the additive (cocycle-trivial)
case.
\end{theorem}

\begin{proof}
By induction on $n$. For $n=1$, both sides equal $v(a_1)$ since the
empty $\omega$-sum is zero. Assume the result for some $n\ge 1$ and
consider $n+1$. By the chain calculus,
$\chain{a_1,\dots,a_{n+1}} = a_1\compose\chain{a_2,\dots,a_{n+1}}$.
Applying the defining equation of $\omega$,
\[
v(\chain{a_1,\dots,a_{n+1}})
= v(a_1) + v(\chain{a_2,\dots,a_{n+1}}) + \omega\bigl(a_1,
  \chain{a_2,\dots,a_{n+1}}\bigr).
\]
Substitute the inductive hypothesis applied to $(a_2,\dots,a_{n+1})$ to
expand $v(\chain{a_2,\dots,a_{n+1}})$ as $\sum_{i=2}^{n+1}v(a_i) +
\sum_{k=2}^{n}\omega(a_k,\chain{a_{k+1},\dots,a_{n+1}})$. Combining
gives exactly the right-hand side at $n+1$. The cocycle identity
$\partial\omega = 0$ guarantees that no other ``rebracketing'' would
yield a different telescoping, but the right-fold is the canonical
one.
\end{proof}

\begin{corollary}[Chain rule for inverses]\label{cor:chain-inverse}
If every $a_i\in\Ideas^{\times}$, then $\chain{a_1,\dots,a_n}\in
\Ideas^{\times}$ with inverse $\chain{a_n^{-1},\dots,a_1^{-1}}$.
\end{corollary}

\begin{proof}
Induction on $n$ via \Cref{prop:inv-group} and the append rule.
\end{proof}

\subsection*{Identity transmission}

A theme of Volume~IV of the Lean development is that the unit $\unit$
is invisible to a chain: inserting $\unit$ anywhere in a chain
preserves the chain.

\begin{theorem}[Identity transmission]\label{thm:id-transmission}
For every list $\mathit{xs}$ of ideas, every position
$1\le k\le \lvert\mathit{xs}\rvert+1$, and every $m\ge 0$, inserting $m$
copies of $\unit$ at position $k$ leaves the chain unchanged. In
particular,
\[
  \chain{\mathit{xs}\!\frown\!(\,\underbrace{\unit,\dots,\unit}_{m})\!\frown\!
         \mathit{ys}}
  \;=\;
  \chain{\mathit{xs}\!\frown\!\mathit{ys}}.
\]
\end{theorem}

\begin{proof}
By the append clause of \Cref{lem:chain-calc}, both sides equal
$\chain{\mathit{xs}}\compose \unit^{(m)}\compose \chain{\mathit{ys}}$
where $\unit^{(m)}$ is the $m$-fold chain of units. By induction on
$m$, $\unit^{(m)} = \unit$. The unit law in the monoid finishes the
calculation. The packaged five-fold version is
\texttt{IdeaTheory.theorem\_7\_2}.
\end{proof}

\subsection*{Filtering and rebracketing}

Two further chain identities will appear in the zoo entries below.

\begin{lemma}[Identity filtering]\label{lem:id-filter}
For any list $\mathit{xs}$ of ideas, the chain
$\chain{\mathit{xs}}$ equals the chain of the sublist obtained by
deleting every occurrence of $\unit$ from $\mathit{xs}$.
\end{lemma}

\begin{proof}
By induction on $\mathit{xs}$. The nil case is trivial. For
$\mathit{xs} = a::\mathit{xs}'$, split on whether $a=\unit$. If
$a=\unit$, both chains equal $\chain{\mathit{xs}'}$ by the unit law.
If $a\ne\unit$, both equal $a\compose\chain{\mathit{xs}'}$ by the
inductive hypothesis applied to $\mathit{xs}'$. The Lean version is
\texttt{IdeaTheory.chain\_filter\_ne\_ident}.
\end{proof}

\begin{lemma}[Block rebracketing]\label{lem:block-reb}
For any partition of a list $\mathit{xs}$ into consecutive blocks
$\mathit{xs} = \mathit{xs}_1 \frown \cdots \frown \mathit{xs}_m$,
\[
  \chain{\mathit{xs}} \;=\; \chain{\chain{\mathit{xs}_1},\dots,
  \chain{\mathit{xs}_m}}.
\]
\end{lemma}

\begin{proof}
By induction on $m$, using the append clause of
\Cref{lem:chain-calc} at each step. This is the algebraic content
of the generalized associative law \cite{Bourbaki1989}: any
parenthesization of a finite product is equivalent.
\end{proof}

\begin{remark}[Operational reading]\label{rem:block-reb}
\Cref{lem:block-reb} is the formal license for a chain-of-chains
analysis: a long transmission can always be summarized as a chain of
its sub-chains, and the value at the top level is identical to the
value computed in one sweep. This is the algebraic substrate of
``levels'' in the cultural-evolution literature \cite{Henrich2015}
and of the multi-scale picture defended by \cite{Sperber1996}.
\end{remark}

\section{Cumulative emergence and the long-chain inequality}
\label{sec:long-chain}

The chain rule (\Cref{thm:chain-rule}) expresses the value of a chain
as the sum of the per-carrier valuations plus a telescoping
$\omega$-correction. When $A$ is ordered (the standard case is
$A=\RR$ with the usual $\le$) and $\omega$ is non-negative, the
correction is monotone and grows with $n$; this is the algebraic
analogue of \emph{cumulative cultural complexity}.

\begin{definition}[Non-negative $\omega$]\label{def:nonneg-omega}
We say the cocycle $\omega\colon\Ideas\times\Ideas\to A$ is
\emph{non-negative} if $A$ is a totally ordered abelian group and
$\omega(a,b)\ge 0$ for all $a,b\in\Ideas$.
\end{definition}

\begin{proposition}[Cumulative emergence is monotone]\label{prop:cum-mon}
If $\omega$ is non-negative, then for every prefix
$(a_1,\dots,a_k)$ of $(a_1,\dots,a_n)$,
\[
v\!\bigl(\chain{a_1,\dots,a_k}\bigr)
- \sum_{i=1}^{k} v(a_i)
\;\le\;
v\!\bigl(\chain{a_1,\dots,a_n}\bigr)
- \sum_{i=1}^{n} v(a_i).
\]
That is, the gap between the chain's value and the sum of its parts
grows monotonically as the chain is extended.
\end{proposition}

\begin{proof}
By \Cref{thm:chain-rule} applied at length $k$ and at length $n$, the
gap at length $k$ is the partial $\omega$-sum
$\sum_{j=1}^{k-1}\omega(a_j,\chain{a_{j+1},\dots,a_k})$, while the gap
at length $n$ is the analogous sum of length $n-1$. The difference is
itself a sum of non-negative $\omega$-terms; non-negativity of
$\omega$ gives the inequality.
\end{proof}

\begin{theorem}[Long-chain inequality]\label{thm:long-chain}
Let $\omega$ be non-negative and let $L\ge 0$ be a uniform lower bound
on $\omega$ over the carriers of interest, i.e.\ assume there is a
constant $L\in A$ with $\omega(a,b)\ge L\ge 0$ for every pair $(a,b)$
that appears as a $(a_k,\chain{a_{k+1},\dots,a_n})$-pair in the
telescoping. Then for every chain of length $n\ge 1$,
\[
v\!\bigl(\chain{a_1,\dots,a_n}\bigr)
\;\ge\;
\sum_{i=1}^{n} v(a_i) + (n-1)\,L.
\]
\end{theorem}

\begin{proof}
Apply \Cref{thm:chain-rule} and bound each of the $n-1$ telescoping
terms below by $L$. The hypothesis is exactly that this bound
applies, and the inequality follows by summation. The case $n=1$ is
vacuous since the sum is empty.
\end{proof}

\begin{remark}[Information-theoretic reading]\label{rem:long-chain}
When $A=\RR$ and $v$ is a complexity valuation (description length,
Kolmogorov complexity, channel cost) the long-chain inequality says
that a sufficiently long transmission produces an output whose
complexity exceeds the sum of the complexities of its inputs by an
amount linear in the length. This is the formal version of
\emph{recombinant growth} \cite{Weitzman1998} and is closely related
to Henrich's \emph{cultural ratchet} (\S\ref{zoo:henrich}).
\end{remark}

\begin{corollary}[No-collapse]\label{cor:no-collapse}
Under the hypotheses of \Cref{thm:long-chain}, if any of $L>0$ holds
strictly, then $\chain{a_1,\dots,a_n}\ne\unit$ for $n$ sufficiently
large, regardless of the choice of carriers.
\end{corollary}

\begin{proof}
By \Cref{thm:long-chain}, $v(\chain{a_1,\dots,a_n})$ exceeds any fixed
bound for $n$ large. Since $v(\unit)=0$ by normalization, the chain
cannot equal $\unit$ once the bound is exceeded.
\end{proof}

\section{Structural equivalence and the universal property}
\label{sec:struct-equiv}

Two idea algebras may be presented very differently --- as a syntactic
word algebra, as a quotient of free vector spaces, as a category of
neural representations --- and yet implement the same structure. The
correct notion of identity between them is structural equivalence,
which we now formalize.

\begin{definition}[Idea algebra morphism]\label{def:morphism}
Let $\IdeaAlg = (\Ideas,\compose,\unit,\rho,\sigma,\nu,\eta,\omega,\deg)$
and $\IdeaAlg' = (\Ideas',\compose',\unit',\rho',\sigma',\nu',\eta',
\omega',\deg')$ be idea algebras. A \emph{morphism} $f\colon\IdeaAlg\to
\IdeaAlg'$ is a function $f\colon\Ideas\to\Ideas'$ such that
\begin{enumerate}
  \item $f(\unit) = \unit'$ and $f(a\compose b) = f(a)\compose' f(b)$;
  \item $\rho'(f(a),f(b)) = \rho(a,b)$ for all $a,b$;
  \item $f(\sigma(a,b))=\sigma'(f(a),f(b))$, and analogously for $\nu$
        and $\eta$;
  \item $\omega'(f(a),f(b)) = \omega(a,b)$;
  \item $\deg'(f(a)) = \phi(\deg(a))$ for some fixed homomorphism
        $\phi\colon G\to G'$ of the grading groups.
\end{enumerate}
A \emph{structural equivalence} is a morphism $f$ that is a bijection
on carriers, with $\phi$ an isomorphism. Structural equivalences form
the morphisms of a category $\mathsf{IdeaAlg}$.
\end{definition}

The bare-monoid version of (1) is the structure
\texttt{IdeaTheory.IdeaHom} in \texttt{Theorems8.lean}: a function
preserving $\unit$ and $\compose$, with the homomorphism law
$f(a\compose b) = f(a)\compose' f(b)$. Definition~\ref{def:morphism}
extends that to all primitives.

\begin{lemma}[Image-of-chain lemma]\label{lem:image-chain}
A morphism $f\colon\IdeaAlg\to\IdeaAlg'$ commutes with chains:
$f(\chain{a_1,\dots,a_n}) = \chain{f(a_1),\dots,f(a_n)}'$, where the
chain on the right uses the operations of $\IdeaAlg'$.
\end{lemma}

\begin{proof}
Induction on $n$, using clause~(1) of \Cref{def:morphism} at the
inductive step. This is the verified \texttt{IdeaTheory.evalWord\_mapWord}.
\end{proof}

We now state the universal property of structural equivalence in two
forms. The first is the \emph{kernel characterization} of the natural
``equality of evaluations'' relation on the free word algebra over a
carrier; the second is a \emph{universal mapping property} for any
structurally equivalent target.

\begin{theorem}[Kernel characterization]\label{thm:kernel}
Let $\Ideas$ be a carrier of an idea algebra and let
$\mathsf{Word}(\Ideas)$ denote the free monoid on $\Ideas$, equipped
with the evaluation map $\mathrm{ev}\colon\mathsf{Word}(\Ideas)\to
\Ideas$ obtained by the chain construction. Define
\[
  w_1 \equivIdea w_2 \;\iff\; \mathrm{ev}(w_1) = \mathrm{ev}(w_2).
\]
Then $\equivIdea$ is the unique relation on $\mathsf{Word}(\Ideas)$
that is (i) an equivalence relation, (ii) a congruence with respect to
\texttt{cons} and append, and (iii) a complete invariant of evaluation:
$w_1\equivIdea w_2$ holds iff $w_1$ and $w_2$ have the same image in
$\Ideas$.
\end{theorem}

\begin{proof}
That $\equivIdea$ is an equivalence is immediate from equality on
$\Ideas$. The congruence claim is verified by combining the cons-rule
$\mathrm{ev}(a::w) = a\compose\mathrm{ev}(w)$ with the append rule
$\mathrm{ev}(w_1\!\frown\!w_2) = \mathrm{ev}(w_1)\compose
\mathrm{ev}(w_2)$ (\Cref{lem:chain-calc}); both
\texttt{IdeaTheory.cons\_congr} and \texttt{IdeaTheory.append\_congr}
record the corresponding lemmas. For uniqueness, suppose $R$ is any
other relation satisfying (i)--(iii). Then for every $w_1,w_2$,
$R(w_1,w_2)$ holds iff $\mathrm{ev}(w_1)=\mathrm{ev}(w_2)$, which is
the definition of $\equivIdea$. The bundled three-clause version is
\texttt{IdeaTheory.theorem\_8\_1}.
\end{proof}

\begin{theorem}[Universal property]\label{thm:universal}
Let $\IdeaAlg$ be an idea algebra. There is a (unique up to unique
isomorphism) idea algebra $\IdeaAlg^{\flat}$ together with a morphism
$\iota\colon\IdeaAlg\to\IdeaAlg^{\flat}$ such that for every idea
algebra $\IdeaAlg'$ and every morphism $f\colon\IdeaAlg\to\IdeaAlg'$
there exists a unique morphism
$\bar f\colon\IdeaAlg^{\flat}\to\IdeaAlg'$ with $\bar f\circ\iota = f$.
Equivalently, the canonical-form construction $w\mapsto[\mathrm{ev}(w)]$
on words is the reflector of $\mathsf{Word}(\Ideas)$ into the
subcategory of structural-equivalence classes.
\end{theorem}

\begin{proof}
Take $\IdeaAlg^{\flat}$ to be the quotient of $\mathsf{Word}(\Ideas)$
by $\equivIdea$, with operations and primitives induced from
$\IdeaAlg$ (well-defined by \Cref{thm:kernel}). The map $\iota$ sends
$a$ to the equivalence class of the singleton word $[a]$. Given
$f\colon\IdeaAlg\to\IdeaAlg'$, define $\bar f$ on a representative
word by composing the letter-by-letter image with the chain
construction in $\IdeaAlg'$, i.e.\ $\bar f([w]) =
\chain{f(a_1),\dots,f(a_n)}'$; \Cref{lem:image-chain} ensures this
descends to the quotient, and uniqueness follows because the singleton
words generate $\IdeaAlg^{\flat}$. The Lean version of the existence
half is \texttt{IdeaTheory.theorem\_8\_2}; the universal-property half
is \texttt{IdeaTheory.theorem\_8\_3}.
\end{proof}

\begin{remark}[\,$\equivIdea$ as a congruence on the original carrier]
\label{rem:cong-carrier}
The same machinery, transported back through evaluation, makes
$\equivIdea$ on $\Ideas$ literally equality (since evaluation is the
identity on singletons). On a non-free carrier the structural
equivalence relation has more bite: it identifies all words with
common evaluation, which by \Cref{lem:chain-calc} encompasses
rebracketing and identity insertion. This is what is verified as
\texttt{IdeaTheory.theorem\_9\_3}: the relation $\mathrm{IdeaEquiv}$
is reflexive, symmetric, transitive, and congruent in both arguments.
\end{remark}

\subsection*{Free algebras and quotient maps}

The reflector $\iota$ of \Cref{thm:universal} has a more concrete
description that will be useful in concrete model-building.

\begin{construction}[Free idea algebra on a set]\label{cons:free}
Let $X$ be a set. The \emph{free idea algebra on $X$}, written
$\IdeaAlg\langle X\rangle$, has as carrier the set
$\mathsf{Word}(X)/\equivIdea$ of equivalence classes of finite words
under structural equivalence with respect to the trivial idea
algebra on $X$ (the one in which $X$ generates a free monoid and all
primitives are extended freely). Composition is concatenation,
$\unit$ is the empty word, $\rho$, $\sigma$, $\nu$, $\eta$, $\omega$,
and $\deg$ are the unique extensions making the universal property
of \Cref{thm:universal} hold.
\end{construction}

\begin{proposition}[Universal mapping property of free algebras]
\label{prop:free-ump}
For every idea algebra $\IdeaAlg'$ and every set-map
$\varphi\colon X\to\Ideas'$ there is a unique morphism
$\bar\varphi\colon\IdeaAlg\langle X\rangle\to\IdeaAlg'$ extending
$\varphi$ in the sense $\bar\varphi(\iota(x)) = \varphi(x)$ for all
$x\in X$.
\end{proposition}

\begin{proof}
Existence: define $\bar\varphi$ on a representative word
$[x_{i_1},\dots,x_{i_n}]$ by chain composition $\chain{\varphi(x_{i_1}),
\dots,\varphi(x_{i_n})}'$; well-definedness on classes is
\Cref{lem:image-chain}; preservation of all primitives uses the
freeness of the construction. Uniqueness: any two such extensions
agree on singleton classes by hypothesis, hence on all classes by
the chain rule and the morphism axioms.
\end{proof}

\section{The headline equivalence theorem}\label{sec:equiv-thm}

We now combine the components into the central result of the chapter.

\begin{theorem}[Structural equivalence theorem]\label{thm:struct-equiv}
Let $\IdeaAlg$ and $\IdeaAlg'$ be idea algebras. The following are
equivalent.
\begin{enumerate}
  \item There exists a structural equivalence $f\colon\IdeaAlg\to
        \IdeaAlg'$.
  \item There exists a graded monoid isomorphism $f\colon\Ideas\to
        \Ideas'$ that intertwines $\rho$, the Aufhebung primitives
        $\sigma,\nu,\eta$, and $\omega$.
  \item For every choice of presentation by generators and chain
        relations, the canonical-form quotients $\IdeaAlg^{\flat}$ and
        $(\IdeaAlg')^{\flat}$ are isomorphic in the category of
        idea algebras, by an isomorphism extending each pair of
        equivalent presentations.
  \item There is a $\rho$-isometric, $\omega$-respecting bijection
        $f\colon\Ideas\to\Ideas'$ such that the conjugation actions of
        $\Ideas^{\times}$ and $(\Ideas')^{\times}$ are intertwined,
        i.e.\ $f(\conj{g}{a}) = \conj{f(g)}{f(a)}$ for all
        $g\in\Ideas^{\times}$ and all $a\in\Ideas$.
\end{enumerate}
\end{theorem}

\begin{proof}
$(1)\Rightarrow(2)$ is immediate from \Cref{def:morphism}: a
structural equivalence is, by definition, a graded monoid isomorphism
satisfying clauses (2)--(5).

$(2)\Rightarrow(1)$. The hypothesis exactly enumerates clauses (1)--(4)
of \Cref{def:morphism}, plus the grading clause (5). Take $\phi$ to
be the induced map on grading groups; it is well-defined and a
homomorphism because $\deg$ is.

$(2)\Leftrightarrow(3)$. By the universal property
(\Cref{thm:universal}), $\IdeaAlg^{\flat}\cong\IdeaAlg$ canonically:
the singleton-word reflector $\iota$ admits a section sending each
class to its evaluation. Hence an isomorphism of canonical forms is
the same as an isomorphism of the original algebras, and the iff
follows.

$(2)\Rightarrow(4)$. By the conjugation theorem
(\Cref{thm:conjugation}), $C_g$ is built out of $\compose$ and inverses
in $\Ideas$, both of which are preserved by $f$ in (2). Hence
$f(C_g(a)) = f(g\compose a\compose g^{-1}) = f(g)\compose'
f(a)\compose' f(g)^{-1} = C_{f(g)}(f(a))$.

$(4)\Rightarrow(2)$. The intertwining property at the special case
$a=\unit$ gives $f(\unit) = \unit'$. The intertwining at $g=$ a
single carrier and $a=$ another carrier, combined with the
$\rho$-isometric and $\omega$-respecting hypotheses, recovers the full
homomorphism property of $f$ on a generating set, and hence on all of
$\Ideas$ by the universal property of \Cref{thm:universal}.
\end{proof}

\begin{corollary}[Conjugation orbits are equivalence invariants]
\label{cor:orbit-inv}
If $f\colon\IdeaAlg\to\IdeaAlg'$ is a structural equivalence, then $f$
restricts to a bijection of conjugacy classes of $\Ideas$ with
conjugacy classes of $\Ideas'$.
\end{corollary}

\begin{proof}
By \Cref{cor:conj-iso} the conjugation action sends conjugacy classes
to conjugacy classes within each algebra. By
\Cref{thm:struct-equiv}(4), $f$ intertwines the actions, hence carries
classes to classes.
\end{proof}

\begin{corollary}[Cohomology class of $\omega$ is an invariant]
\label{cor:omega-cohomology}
If $\IdeaAlg\cong\IdeaAlg'$ as idea algebras, then the cohomology
classes $[\omega]\in H^2(\Ideas;A)$ and $[\omega']\in H^2(\Ideas';A)$
agree under the induced isomorphism on cohomology. In particular, the
emergence-cohomology class is a complete obstruction to additivity:
two algebras can be structurally equivalent only if their emergence
defects are cohomologous.
\end{corollary}

\begin{proof}
\Cref{thm:struct-equiv}(2) yields $\omega'(f(a),f(b)) = \omega(a,b)$
on the nose; this is a stronger statement than mere cohomologousness,
but it certainly implies it. The induced map on $H^2$ is functorial in
the carrier monoid, by the standard machinery of group cohomology
\cite{Weibel1994,Brown1982}.
\end{proof}

\begin{corollary}[Grading is a complete first-order invariant]
\label{cor:deg-invariant}
If $\IdeaAlg$ and $\IdeaAlg'$ are structurally equivalent, then the
multisets $\{\deg(a) : a\in\Ideas\}$ and $\{\deg'(a') : a'\in\Ideas'\}$
agree as multisets in their respective grading groups (which are
themselves isomorphic).
\end{corollary}

\begin{proof}
A structural equivalence is a bijection that preserves $\deg$ up to
the isomorphism $\phi\colon G\to G'$.
\end{proof}

\section{The zoo}\label{sec:zoo-ch6}

We close with seven entries from the zoo of ideas about ideas, each
read through the formalism of conjugation, transmission, and
structural equivalence. The pattern is the standard one:
original claim, formal restatement, status under the theorems above,
and brief commentary.

\subsection{Bourdieu --- habitus as conjugation by class position}
\label{zoo:bourdieu}

\begin{zooentry*}
\thinker{Pierre Bourdieu}, \emph{Distinction} \cite{Bourdieu1984} and
the wider field-and-habitus programme, treats the social agent as
carrying a system of dispositions --- the habitus --- that is
\emph{class-positional}: the same gesture, taste, or judgement
acquires different significance depending on the position from which
the agent enacts it. A bourgeois reading of \emph{The Magic Mountain}
is not the same idea as a working-class reading of the same novel,
even though the propositional content is identical.
\formalclaim Each class position $g$ is modelled as an invertible idea
$g\in\Ideas^{\times}$ encoding the social-cognitive frame, and the
class-positioned act is $\conj{g}{a} = g\compose a\compose g^{-1}$.
The habitus is the orbit map $a\mapsto \{\conj{g}{a} : g\in
G_{\mathrm{soc}}\}$, where $G_{\mathrm{soc}}\subseteq\Ideas^{\times}$
is the subgroup of socially realized invertibles.
\formalstatus
True \emph{up to} the conjugation theorem (\Cref{thm:conjugation}):
$\rho$, $\omega$, $\deg$ are invariant along class-positional orbits,
which is what allows Bourdieu to insist that ``misrecognition'' of
class effects is structurally invisible to the agents involved. The
formalism predicts in addition that any quantity \emph{not} expressible
in $\rho$, $\sigma$, $\nu$, $\eta$, $\omega$, or $\deg$ may
discriminate between classes; Bourdieu's empirical work on body
hexis, on which he draws elsewhere \cite{Granovetter1973} for
network-mediated transmission, supplies exactly such non-invariants.
\litcomment
The fit is unusually clean. Bourdieu's persistent claim that ``the
field structures the habitus'' becomes the algebraic claim that
$g\mapsto C_g$ is a homomorphism. The persistent worry that
``positionality'' might be a trivial relabeling --- and hence not
explanatory --- is dispelled by \Cref{cor:orbit-inv}: an orbit of
$C_g$ identifies socially equivalent acts as algebraically equivalent
ones, but the orbit itself is a non-trivial structural invariant
under structural equivalence.
\end{zooentry*}

\subsection{Skinner --- context as conjugation}\label{zoo:skinner}

\begin{zooentry*}
\thinker{Quentin Skinner}, \emph{Meaning and understanding in the
history of ideas} \cite{Skinner1969}, argues against the practice of
reading historical texts as direct expressions of timeless doctrines.
The meaning of a Hobbesian utterance about ``commonwealth'' is not
the meaning of the same English word for a 21st-century reader; it is
constituted by the surrounding intellectual context --- what
arguments could have been made, what conventions were available,
which interlocutors were targeted.
\formalclaim Each intellectual context $g$ is an invertible operator
on the ideational content; the contextualized utterance is
$\conj{g}{a}$. Skinner's methodology demands that the historian
recover $g$ before evaluating $a$.
\formalstatus
The conjugation theorem makes this precise: $\rho(\conj{g}{a},
\conj{g}{b}) = \rho(a,b)$ shows that \emph{relative} resonances
between contextualized utterances reproduce the relative resonances
of the bare ones. Hence Skinner's claim that ``meaning is preserved
modulo context'' is exactly Idea Theory's invariance of $\rho$ under
inner automorphisms. By contrast, $v(\conj{g}{a})$ for a non-trivial
valuation $v$ generally differs from $v(a)$ by a coboundary
\cite{Brown1982}, which models the historiographical fact that
quantitative scoring (rhetorical force, persuasive uptake) varies
across contexts.
\litcomment
Skinner's own critique of ``perennial questions'' \cite{Skinner1969}
becomes the rejection of the assumption that the orbit of $a$ under
$C_g$ is a singleton. Idea Theory grants Skinner the formal
machinery to say so without committing him to relativism: the orbit
exists as an algebraic object, and is fully invariant under
structural equivalence.
\end{zooentry*}

\subsection{Foucault --- épistémè as a conjugation orbit}\label{zoo:foucault}

\begin{zooentry*}
\thinker{Michel Foucault}, \emph{L'Arch\'eologie du savoir}
\cite{Foucault1969}, proposes that knowledge in any historical period
is organized by a deep structural form --- the \emph{épistémè} ---
that determines what can be said, thought, and counted as a
legitimate object of knowledge. An épistémè is not a doctrine but a
configuration: it conjugates every discursive act it touches.
\formalclaim An épistémè is modelled as a (large) subgroup
$E\subseteq\Ideas^{\times}$. The discursive content of a period is
the orbit set $\Ideas/E = \{\,E\!\cdot\! a : a\in\Ideas\,\}$, where
the action is by inner automorphism. A \emph{discursive rupture} is a
period in which the appropriate $E$ changes; this is exactly the
absence of an inner $g$ that would conjugate the old discursive
configuration into the new.
\formalstatus
True \emph{conditionally} on \Cref{thm:conjugation} and on
\Cref{thm:struct-equiv}(4). Within a single épistémè, the conjugation
theorem guarantees that all primitives are preserved along orbits, so
two utterances in the same orbit are algebraically indistinguishable
--- the formal sense of Foucault's ``regularity of statements''
\cite{Foucault1969}. Across épistémès, the ``rupture'' phenomenon is
the failure of any structural equivalence
(\Cref{thm:struct-equiv}(1)) to extend the within-period
conjugations: there simply is no graded monoid isomorphism between
classical-age and modern-age $\Ideas$.
\litcomment
This is the sharpest of the seven fits in this chapter. Foucault's
notorious refusal of any continuous-history account is the algebraic
fact that not every cross-period bijection is a structural
equivalence, and his equally insistent claim that within a period
\emph{everything} is reorganized along the orbit lines is the
conjugation theorem.
\end{zooentry*}

\subsection{Latour --- translation chains}\label{zoo:latour}

\begin{zooentry*}
\thinker{Bruno Latour}, \emph{Reassembling the Social}
\cite{Latour2005} and the broader Actor-Network Theory programme,
analyzes scientific and social facts as products of long chains of
\emph{translations}: every step --- a measurement, an inscription, a
publication, an institutional uptake --- transforms the content
slightly while purporting to transmit it. The gap between the
original and the final product is not noise but the structuring
phenomenon itself.
\formalclaim Each translation step is a carrier $a_i\in\Ideas$, and
the resulting fact is the chain $\chain{a_1,\dots,a_n}$. The
``cost of translation'' for a step is the local emergence
$\omega(a_k,\chain{a_{k+1},\dots,a_n})$ of \Cref{thm:chain-rule};
the cumulative cost is the long-chain sum.
\formalstatus
True under the chain rule (\Cref{thm:chain-rule}) and the long-chain
inequality (\Cref{thm:long-chain}). The non-additivity of $v$ along
the chain is exactly Latour's insistence that the final fact is
\emph{not} the sum of its mediators. The non-negativity assumption on
$\omega$ corresponds to ANT's empirical claim that translation is
\emph{costly}: each carrier adds something, and successive additions
accumulate. Latour's own examples (the Pasteurization of France, the
inscription of soil-pH measurements at Boa Vista) trace exactly the
$n-1$ telescoping terms.
\litcomment
The fit is literal: ``transmission chain'' and ``translation chain''
denote the same algebraic object, and the cost-of-translation
intuition is the algebraic content of the cocycle. The corollary of
no-collapse (\Cref{cor:no-collapse}) gives a Latourian prediction:
sufficiently long networks always produce non-trivial output, even
from trivial inputs --- which is exactly the assertion that
``society'' is constituted, not given.
\end{zooentry*}

\subsection{Rogers --- diffusion of innovations}\label{zoo:rogers}

\begin{zooentry*}
\thinker{Everett Rogers}, \emph{Diffusion of Innovations}
\cite{Rogers2003}, models the spread of an innovation across a
social system as a sequence of adoption decisions, each modulated by
the adopter's awareness, persuasion, decision, implementation, and
confirmation phases. Long adoption chains exhibit increasing
\emph{resistance}: late adopters face higher cumulative friction
than early adopters.
\formalclaim Each adopter or adoption step is an idea $a_i$; the
total ``adoption work'' done up to time $n$ is
$v(\chain{a_1,\dots,a_n})$. The resistance encountered by step $k$ is
the local cocycle term $\omega(a_k,\chain{a_{k+1},\dots,a_n})$, and
the cumulative resistance is the second sum of \Cref{thm:chain-rule}.
\formalstatus
True under \Cref{thm:long-chain}: if the per-step resistance is
bounded below by some $L>0$, then the total resistance grows
linearly in chain length, exactly the empirical S-curve plateau
documented by Rogers \cite{Rogers2003}. The strict version
\Cref{cor:no-collapse} says that long-chain adoptions cannot reduce
to the trivial idea, which corresponds to Rogers's observation that
mature diffusion produces social structure that does not vanish on
discontinuation of the original innovation.
\litcomment
Rogers's five-stage model maps less directly onto the formalism than
Latour's translation, because the stages are heterogeneous; what
survives translation is the long-chain inequality, which in
sociological terms is the well-attested ``cumulative friction''
phenomenon. The bias toward early adopters --- a structural
prediction of the model --- corresponds in Idea Theory to the
position-dependence of $\omega(a_k, \chain{a_{k+1},\dots,a_n})$ on
the depth of the suffix.
\end{zooentry*}

\subsection{Henrich --- the cultural ratchet}\label{zoo:henrich}

\begin{zooentry*}
\thinker{Joseph Henrich}, \emph{The Secret of Our Success}
\cite{Henrich2015}, argues that human cognitive distinctiveness
arises from a \emph{cultural ratchet}: the cumulative, irreversible
build-up of know-how across generations of social transmission. Each
generation inherits a baseline and adds incrementally; the result is
artefacts and skills no individual could have invented.
\formalclaim Each generation is a chain of carriers, and the
cross-generational accumulation is the iterated chain
$\chain{\chain{a_{1,1},\dots,a_{1,n_1}},\dots,\chain{a_{k,1},\dots,
a_{k,n_k}}}$. The ratchet's monotonicity is
\Cref{prop:cum-mon}: the cumulative emergence is non-decreasing in
chain extension when $\omega\ge 0$.
\formalstatus
True under \Cref{prop:cum-mon} and \Cref{thm:long-chain}. The
\emph{irreversibility} that Henrich emphasizes is not literal
algebraic irreversibility (the chain rule gives a closed-form
inverse for invertible carriers, \Cref{cor:chain-inverse}) but the
\emph{positive monotone growth} of cumulative $\omega$. The
no-collapse corollary (\Cref{cor:no-collapse}) is the algebraic form
of Henrich's experimental observation that long chains of social
learners exceed any individual's reasoning capacity \cite{Henrich2015}.
\litcomment
Henrich's empirical case for cultural over individual cognition rests
on chain-length effects; Idea Theory predicts these effects must be
of the form $(n-1)L$ to leading order, with $L$ a per-step cocycle
floor. The formalism gives a sharper hypothesis than the verbal
``ratchet'' --- it predicts that the slope, not just the sign, of the
cumulative-emergence curve is structurally invariant under
\Cref{thm:struct-equiv}.
\end{zooentry*}

\subsection{Kuhn --- paradigm as structural equivalence class}
\label{zoo:kuhn}

\begin{zooentry*}
\thinker{Thomas Kuhn}, \emph{The Structure of Scientific Revolutions}
\cite{Kuhn1962}, distinguishes ``normal science'' --- routine
puzzle-solving within a shared paradigm --- from ``revolutionary
science'', in which one paradigm replaces another, and the two are in
some strong sense \emph{incommensurable}: terms like ``mass'' refer
differently across the divide.
\formalclaim A paradigm is a structural-equivalence class of idea
algebras: the collection of all algebras isomorphic to a fixed
$\IdeaAlg$ in the sense of \Cref{def:morphism}. Normal science is
work within one such class; a revolution is the absence of any
structural equivalence connecting the pre- and post-revolutionary
algebras.
\formalstatus
True under \Cref{thm:struct-equiv}. Within a class, every primitive
($\rho$, $\sigma$, $\nu$, $\eta$, $\omega$, $\deg$) is preserved
under translation between presentations, and the conjugation action
is intertwined; this is the algebraic version of Kuhn's
``paradigm-internal articulation''. Revolutions are diagnosed by the
failure of clause~(2) of \Cref{thm:struct-equiv}: there is no
$\rho$-isometric, $\omega$-respecting, graded monoid isomorphism
extending any of the available translations. Kuhn's famous claim that
revolutions cannot be resolved by argument internal to either
paradigm is then a theorem: argument internal to $\IdeaAlg$ moves
within its conjugacy structure, and conjugation is structurally
invariant (\Cref{cor:orbit-inv}).
\formalclaim
\litcomment
Two predictions of the formalism that go beyond Kuhn. First, by
\Cref{cor:omega-cohomology}, the obstruction to a structural
equivalence between paradigms is concentrated in the cohomology
class of $\omega$: paradigms can fail to be structurally equivalent
only if their emergence cocycles are non-cohomologous. Second, by
\Cref{cor:deg-invariant}, the multiset of grading values is a
fingerprint of a paradigm; ``hidden'' structural equivalences --- not
detected by surface translations --- can be ruled out simply by
counting graded levels. Both refinements give Kuhn's broadly literary
thesis a quantitative bite that subsequent commentators
\cite{Lovejoy1936,Koselleck2002} have demanded.
\end{zooentry*}

\section*{Summary of the chapter}

\section{Coda: methodological consequences}\label{sec:coda-ch6}

The three devices of this chapter --- conjugation, transmission, and
structural equivalence --- act in concert in every later chapter, and
it is worth pausing to draw out the methodological lessons they
license. Each is briefly stated and then situated.

\paragraph{Frame-relative invariance.} The conjugation theorem has the
methodological force of a Galilean invariance principle: any property
expressible in $\rho$, $\sigma$, $\nu$, $\eta$, $\omega$, or $\deg$ is
``frame-independent'' in the precise sense of being constant along
inner-automorphism orbits. Empirical disputes about the meaning of an
idea that turn on the choice of frame --- ``but in our paradigm $a$
means $a'$'' --- are resolved by exhibiting an invertible $g$ with
$\conj{g}{a} = a'$. The dispute is then revealed to be merely about
presentation, since structural equivalence (\Cref{thm:struct-equiv})
will identify the two sides. Conversely, disputes that resist any
such inner $g$ are diagnosed as genuine: they signal a non-trivial
quotient of $\Ideas^{\times}$ acting on $\Ideas$.

\paragraph{Chain accounting.} The chain rule of \Cref{thm:chain-rule}
turns the otherwise-vague injunction ``track the transmission'' into
an explicit telescoping. Whenever an analyst is tempted to write
``content is roughly preserved across stages'', the chain rule asks
which valuation $v$, which cocycle $\omega$, and which decomposition
of the chain are intended. The long-chain inequality
(\Cref{thm:long-chain}) supplies a non-trivial \emph{quantitative}
prediction --- linear cumulative emergence --- that an empirical
analyst can either confirm or refute, and that the formalism cannot
be made to evade by re-bracketing (\Cref{lem:block-reb}) or by
identity insertion (\Cref{thm:id-transmission}).

\paragraph{Equivalence as a benchmark, not a definition.}
\Cref{thm:struct-equiv} gives four mutually equivalent
characterizations of when two idea algebras are ``the same''. None is
\emph{the} definition; each highlights a different working face of
equivalence. The morphism formulation (clause~1) is best for
constructive arguments; the algebraic-data formulation (clause~2) is
best for case-by-case verification; the canonical-form formulation
(clause~3) is best for computational implementations; and the
intertwining-conjugation formulation (clause~4) is best when the
interest is in symmetry rather than in structure tout court. We will
draw on each of the four in different chapters.

\paragraph{Cohomological obstructions.} \Cref{cor:omega-cohomology}
isolates a single, computable obstruction to structural equivalence:
the cohomology class of $\omega$. This sharpens the otherwise
qualitative discourse of ``incommensurability''
\cite{Kuhn1962,Foucault1969,Bedau1997} by giving the field a concrete
invariant to compute. In the cases where $H^2$ is finite-dimensional
--- the typical case for finitely generated carriers
\cite{Brown1982,Weibel1994} --- the diagnosis whether two
presentations are or are not structurally equivalent reduces to a
finite cohomological calculation.

\paragraph{Looking ahead.} Chapter~\ref{ch:grading} will exploit the
grading invariance of \Cref{thm:conjugation}(5) to develop the theory
of \emph{degree filtrations}: the chain $\deg_0\subseteq\deg_1\subseteq
\cdots$ of subalgebras at successive grading levels. The conjugation
action restricts to each level, and the chain rule respects degree by
the same argument as in the proof of \Cref{thm:conjugation}(5).
Chapter~\ref{ch:cogsci} then specializes the transmission-chain
formalism to cognitive carriers (working-memory traces, articulatory
gestures, neural activation patterns), where the long-chain
inequality recovers a battery of empirical findings on serial recall
and cumulative-cultural complexity \cite{Henrich2015,Sperber1996}.
Chapter~\ref{ch:econ_game} will treat structural equivalence as an
\emph{equilibrium} concept in evolutionary game-theoretic models of
ideational competition \cite{MaynardSmith1982}: two strategies are
indistinguishable in equilibrium iff they are structurally
equivalent in the sense of \Cref{thm:struct-equiv}.

\section*{Summary of the chapter}

We isolated $\Ideas^{\times}$, proved it is a group, and showed that
its inner action on $\Ideas$ preserves every primitive of the
idea-algebra structure (\Cref{thm:conjugation}). We folded finite
sequences of carriers into transmission chains and proved a chain rule
for $\omega$ (\Cref{thm:chain-rule}), a long-chain inequality
(\Cref{thm:long-chain}), and an identity-transmission theorem
(\Cref{thm:id-transmission}). We defined morphisms of idea algebras and
proved the universal property (\Cref{thm:universal}) and the
structural-equivalence theorem (\Cref{thm:struct-equiv}), with
corollaries on conjugacy invariance, cohomology, and grading. The zoo
entries showed that habitus, contextualism, épistémè, translation, the
diffusion of innovations, the cultural ratchet, and Kuhnian paradigms
admit clean formal restatements in this language. The next chapter
takes up the grading $\deg$ in its own right.
